% ======================================================================
%  Section 4.7.2 --- Convergence under occlusion.  New text.
%  The section currently contains Figures 4.6-4.9 and Table 4.2 and no prose.
%
%  MEASURED VALUES (iteration 600, converted from the console):
%
%    variant     |dt|        |dw|       misalignment   |Vc| (iter 300)
%    1 none      1291 um     1.6953 deg   25.8 px      1.77e-6
%    2 Tukey     2.91 um     0.0164 deg    0.25 px     3.27e-5
%    3 Hermite   3.12 um     0.0152 deg    0.23 px     7.59e-6
%
%    |e|^2 at iteration 300: 1.551e10, 1.573e10, 1.573e10
%    %%% NEEDED: the same three values at iteration 600. The argument in the
%    %%% third paragraph depends on Variant 1 still having the SMALLEST
%    %%% feature error at 600, as it does at 300. Verify before keeping it.
%
%  WHAT THIS SECTION MUST NOT CLAIM. Variants 2 and 3 differ by 7.6% in
%  rotation and 7.2% in translation, in opposite directions. They are
%  indistinguishable here. Section 4.6.1 calls their comparison the decisive
%  one of the chapter, so the tie has to be reported as a tie.
%
%  WHAT IT MUST NOT DO EITHER. Do not compare these numbers with Table 4.1's.
%  The occluded runs reach micrometre accuracy and the nominal runs do not,
%  which would suggest that occlusion improves positioning. The cause is the
%  gain schedule: lambda is proportional to the intensity error, the occluder
%  puts a permanent floor under it (14.8 grey levels RMS against 2.4 in the
%  nominal case), so lambda stays near 2.2 instead of decaying to 0.04 and
%  the controller keeps driving. Within the occluded case all three variants
%  share lambda to within 6%, so the three-way comparison is clean.
% ======================================================================

\subsection{Convergence under occlusion}
\label{sec:ablation-results}

The target is now occluded by a disc covering $7.4\%$ of the measurements,
present in the images acquired during servoing but absent from the reference
image, so that the corrupted measurements are genuine outliers with respect to
the model rather than a change of goal. Figure~\ref{fig:occ-setup} shows the
configuration. Everything else is as in Section~\ref{sec:nominal}: same initial
and desired poses, same bank, same gain schedule, same iteration cap.

\paragraph{The unweighted scheme converges to the wrong pose.}
Variant~1 terminates $1.29$~mm and $1.70^{\circ}$ from the desired pose, a
misalignment of some $26$~pixels at $f=870$~px. This is not slow convergence:
its commanded velocity has fallen to $1.8\times10^{-6}\,\mathrm{m\,s^{-1}}$,
\emph{smaller} than that of either weighted variant, and
Figure~\ref{fig:occ-v1}(b) shows the pose error flat over the last three
hundred iterations. The scheme has converged, and it has converged somewhere
else. The occluded block biases the descent direction exactly as
Section~\ref{sec:least-squares} predicts of a quadratic cost, and the
minimiser of the corrupted objective is a quarter of a degree and a millimetre
away from the minimiser of the uncorrupted one.

\paragraph{Both weighted variants recover the pose.}
Variants~2 and~3 terminate $2.9$ and $3.1$~micrometres from the desired pose,
with rotational errors of $0.016^{\circ}$ and $0.015^{\circ}$. Relative to the
baseline this is an improvement of a factor of roughly $400$ in translation
and $100$ in rotation, and it takes the residual misalignment from $26$~pixels
to a quarter of a pixel. The weighting does what Section~\ref{sec:weighted-law}
claims for it: the corrupted measurements contribute neither to the descent
direction nor to the curvature estimate, and the pose is recovered from the
remainder.

\paragraph{The two weightings are indistinguishable at this occlusion level.}
Variant~3 attains the smaller rotational error, $0.0152^{\circ}$ against
$0.0164^{\circ}$, and the larger translational error, $3.12$ against
$2.91\,\mu$m --- differences of $7.6\%$ and $7.2\%$ respectively, in opposite
directions. On a single run per variant no ordering can be read from
differences of that size. The comparison identified in
Section~\ref{sec:ablation-variants} as the decisive one of the chapter
therefore returns a tie: on this configuration the structural information does
not measurably improve on generic robust estimation.

What it does change is which measurements are rejected and when, and that is
not visible in the pose error at all. Section~\ref{sec:weight-maps} examines
it directly and finds that the two mechanisms differ substantially even where
their outcomes do not.

Two qualifications bound this result rather than weaken it. The occluder here
is a single compact disc of uniform intensity, which is the case the structure
measure of~\eqref{eq:structG} is best suited to and also the case a
residual-only estimator finds easiest, since the corrupted residuals are large
and mutually consistent. The configurations in which the two might be expected
to separate --- a fragmented occluder, whose perimeter-to-area ratio
Section~\ref{sec:mad} identifies as the relevant quantity, or an occluded
fraction approaching the breakdown point --- are not tested here and are
identified as future work in Section~\ref{sec:future}.

\paragraph{The feature error is not a proxy for accuracy.}
%%% VERIFY at iteration 600 before keeping this paragraph; the values quoted
%%% are from iteration 300.
The three variants terminate at feature errors of $1.551\times10^{10}$,
$1.573\times10^{10}$ and $1.573\times10^{10}$. Variant~1 therefore attains the
\emph{smallest} feature error of the three while positioning the camera a
hundred times less accurately. This is the clearest demonstration in the
chapter of a point made in passing in Section~\ref{sec:nominal}: the
unweighted scheme minimises the corrupted objective more successfully than the
weighted ones do, and the minimum of the corrupted objective is in the wrong
place. Any assessment of these variants by their residual would rank them
exactly backwards.

\paragraph{A caution on comparing the two cases.}
The weighted variants reach micrometre accuracy here and sub-millimetre
accuracy in the nominal case of Section~\ref{sec:nominal}, which would seem to
say that occlusion improves positioning. It does not. The cause is the gain
schedule. By~\eqref{eq:lambda} the adaptive gain is proportional to the
intensity error, and the occluder places a permanent floor under that error:
it settles at $14.8$ grey levels root-mean-square against $2.4$ in the nominal
case, because no pose can cancel a residual that arises from a difference in
scene content. The gain accordingly remains near $\lambda=2.2$ instead of
decaying to $\lambda=0.04$, and the controller continues to drive the camera
long after it has stopped in the nominal case. Nothing about the occlusion
assists the positioning; the scheme simply never stops pushing. Within the
occluded case the three variants share $\lambda$ to within six per cent, so
the comparison made above is unaffected; across the two cases the numbers are
not comparable and should not be read as such.

% ======================================================================
%  LABELS ASSUMED --- rename to match your own:
%    fig:occ-setup            Figure 4.6
%    fig:occ-v1               Figure 4.7
%    sec:least-squares        4.2.1
%    sec:weighted-law         4.2.2
%    sec:ablation-variants    4.6.1
%    sec:weight-maps          4.7.3
%    sec:nominal              4.7.1
%    sec:mad                  4.3.2
%    sec:future               future work
% ======================================================================
